% 02-difficultes-mecanique.tex — État de l'art 1/5 : difficultés des étudiants
% en mécanique de première année (≈ 6 pages).

\chapter{Apprendre la mécanique en première année}
\label{chap:difficultes}

\section{La mécanique dans le cursus d'ingénieur civil}

\begin{arediger}{1 page}
Place du cours de Physique~1 (EPL1201) : contenus, formulaire officiel
\parencite{formulaire_epl1201}, notations imposées (\(u_0\), \(v_0\), MRU/MRUA,
\(a_r = -r \omega^2\)). Pourquoi ces notations comptent pour les
simulations (cohérence avec le cours, voir section~\ref{sec:notation}).
\end{arediger}

\section{Conceptions alternatives sur la force et le mouvement}

\begin{arediger}{2 pages}
Littérature sur les préconceptions : force dans le sens du mouvement,
\enquote{force centrifuge}, confusion vitesse/accélération, poids et
réaction normale toujours égaux. Le Force Concept Inventory
\parencite{hestenes1992fci} et ses résultats typiques en début d'université.
Lien direct avec nos exercices : réaction normale variable dans le looping,
tension qui s'annule au sommet.
\end{arediger}

\section{Le mouvement circulaire et la réaction normale}

\begin{arediger}{1,5 page}
Difficultés spécifiques : accélération centripète sans \enquote{force
centripète} supplémentaire, rôle de \(N\) et de \(T\), condition \(N \geq 0\).
Le looping comme situation classique (manuels, attractions ;
\textcite{pendrill2005loops} sur la forme des loopings réels). Présenter la
figure~\ref{fig:etat-normale} : \(N/(m g)\) le long du cercle pour trois
hauteurs.
\end{arediger}

\begin{figure}[ht]
  \centering
  \begin{tikzpicture}
    \begin{axis}[width=0.8\linewidth, height=5.4cm, xmin=0, xmax=360, xtick={0,90,180,270,360},
      xlabel={angle depuis le bas du looping \(\theta\) (°)}, ylabel={\(N / (m g)\)},
      grid=major, grid style={gray!20}, legend pos=north east, legend cell align=left,
      every axis plot/.append style={thick}]
      \addplot[accent] table[x=theta, y=h30, col sep=comma] {figures/donnees/normale_looping.csv};
      \addplot[normale] table[x=theta, y=h25, col sep=comma] {figures/donnees/normale_looping.csv};
      \addplot[poids] table[x=theta, y=h20, col sep=comma] {figures/donnees/normale_looping.csv};
      \addplot[black, dashed, thin] coordinates {(0,0) (360,0)};
      \legend{\(h = 3R\), \(h = 5R/2\) (limite), \(h = 2R\)}
    \end{axis}
  \end{tikzpicture}
  \caption{Réaction normale réduite \(N/(m g) = 2 (h - y)/R + \cos(\theta)\) le long
    du looping, avec \(y = R (1 - \cos(\theta))\). Pour \(h = 5R/2\), \(N\) s'annule
    exactement au sommet ; pour \(h = 2R\), elle devient négative : le wagon quitte
    le rail. Données générées par le code de l'application (\fichier{outils/donnees.mjs}).}
  \label{fig:etat-normale}
\end{figure}

\section{L'énergie mécanique comme outil de raisonnement}

\begin{arediger}{1,5 page}
Conservation de l'énergie, choix d'une référence, passage énergie
\(\rightarrow\) vitesse \(\rightarrow\) force. Difficulté à combiner deux
principes (énergie et deuxième loi de Newton) dans un même problème : c'est
exactement ce qu'exige la condition \(h_{\min} = 5R/2\).
\end{arediger}
